\documentclass[11pt]{article}
\usepackage[margin=1in]{geometry}
\usepackage{amsmath,amssymb,amsthm,mathtools,booktabs}
\usepackage[colorlinks=true,linkcolor=blue,urlcolor=blue]{hyperref}
\newcommand{\Tr}{\operatorname{Tr}}
\newcommand{\tr}{\operatorname{tr}}
\newcommand{\ee}{\varepsilon}
\newtheorem{theorem}{Theorem}
\newtheorem{lemma}{Lemma}
\title{The thermal-purity recursion in \(2+1\) dimensions}
\author{}
\date{October 8, 2026}
\begin{document}
\maketitle

\begin{abstract}
	For a frustration-free Hamiltonian with exact ground energy zero, the
	natural recursion variable is the partition-function excess \(R=Z-1\).
	Temporal doubling squares this excess from above. Two scalar spatial
	gluing inequalities then give
	\(R'\le(1+R^2)^4-1\le5R^2\) for \(R\le0.01\).
	A single finite-volume seed yields a uniform gap on dyadic tori;
	spatial moment interpolation extends it to intermediate periods.
	We also retain the general three-purity cube recursion, derive an
	error-tolerant gluing update, and give an exact local obstruction to
	the standard time-reversal reflection-positive construction for AKLT.
	The scalar spatial gluing inequalities remain unproved here for AKLT.
	For a separate exact projection filter, we certify quadratic spatial
	concentration at transverse width two and depth one despite nonreal
	transfer eigenvalues. Extending this estimate as width and depth grow
	remains open.
\end{abstract}

\section{Quadratic Excess Recursion for Frustration-Free Systems}

Suppose \(H_{n,m}\ge0\) has at least one zero-energy ground state for
each periodic rectangle considered. Set
\[
	Z(n,m,\beta)=\Tr e^{-\beta H_{n,m}},\qquad
	R(n,m,\beta)=Z(n,m,\beta)-1.
\]
One zero-energy eigenvector contributes exactly one to \(Z\). If the
remaining energies, including any additional zeros, are \(E_k\), then
\begin{equation}
	R(n,m,2\beta)=\sum_{k\ge1}e^{-2\beta E_k}
	\le\left(\sum_{k\ge1}e^{-\beta E_k}\right)^2
	=R(n,m,\beta)^2.
	\label{eq:ff-time}
\end{equation}
This does not assume uniqueness. A bound \(R<1\) subsequently proves
uniqueness, since any additional zero-energy state contributes one.

The spatial input is only the pair of scalar inequalities
\begin{align}
	Z(2n,m,b)  & \le Z(n,m,b)^2,\label{eq:ff-x}  \\
	Z(2n,2m,b) & \le Z(2n,m,b)^2.\label{eq:ff-y}
\end{align}
For one recursion step these are needed at \(b=2\beta\).
An exact self-adjoint spatial-transfer representation implies them,
but is a stronger assumption: no spatial squaring lemma is used here.

\begin{theorem}[Frustration-free \(2+1\)D recursion]
	If \eqref{eq:ff-x}--\eqref{eq:ff-y} hold at \(b=2\beta\), then
	\begin{equation}
		\boxed{R(2n,2m,2\beta)
			\le\bigl[1+R(n,m,\beta)^2\bigr]^4-1.}
		\label{eq:ff-update}
	\end{equation}
	In particular, for \(0\le R\le0.01\),
	\begin{equation}
		\boxed{R'\le4R^2+6R^4+4R^6+R^8\le5R^2.}
		\label{eq:ff-quadratic}
	\end{equation}
\end{theorem}

\begin{proof}
	First double imaginary time using \eqref{eq:ff-time}; then double
	the two spatial periods using \eqref{eq:ff-x}--\eqref{eq:ff-y}:
	\[
		Z(2n,2m,2\beta)\le Z(n,m,2\beta)^4
		\le\bigl[1+R(n,m,\beta)^2\bigr]^4.
	\]
	Subtract one. For \(R\le0.01\), the coefficient of \(R^2\) is at
	most \(4+6(0.01)^2+4(0.01)^4+(0.01)^6<5\).
\end{proof}

The order matters for the bound: temporal doubling first gives the
leading coefficient \(4\). Spatial doubling first would give the
weaker \(R'\le[(1+R)^4-1]^2\), with leading coefficient \(16\).
In \(d\) spatial dimensions the same argument gives
\(R'\le(1+R^2)^{2^d}-1\), with leading coefficient \(2^d\).

\paragraph{One finite seed.}
Assume the gluing inequalities at every subsequent dyadic step, and
certify only
\begin{equation}
	Z(n_0,m_0,\beta_0)\le1.01.
	\label{eq:ff-seed}
\end{equation}
With \((n_j,m_j,\beta_j)=2^j(n_0,m_0,\beta_0)\), induction gives
\begin{equation}
	R(n_j,m_j,\beta_j)\le r_j:=\frac15\,20^{-2^j},
	\qquad r_{j+1}=5r_j^2.
	\label{eq:ff-scales}
\end{equation}
The ground state is unique on each such torus and
\(e^{-\beta_j\Delta_{n_j,m_j}}\le R(n_j,m_j,\beta_j)\).
Thus \(\Delta_{n_j,m_j}>(\log20)/\beta_0\), uniformly in \(j\).
There is no need to reconstruct the gap from a normalized purity.

If a finite seed Hamiltonian has a unique zero-energy ground state,
Hilbert-space dimension \(D\), and a certified finite-volume gap
\(\gamma_*>0\), then
\[
	R(n_0,m_0,\beta)\le(D-1)e^{-\beta\gamma_*}.
\]
Consequently, \(\beta_0=\log[100(D-1)]/\gamma_*\) suffices for
\eqref{eq:ff-seed}. Under the spatial assumptions, this yields the
explicit, possibly small, lower bound
\begin{equation}
	\Delta_{n_j,m_j}>
	\frac{\gamma_*\log20}{\log[100(D-1)]}.
	\label{eq:ff-finite-gap}
\end{equation}
Existence of a finite-volume seed is therefore automatic from finite
dimensionality and uniqueness. The substantive uniform input is the
spatial gluing property, not the existence of three low-temperature
purity certificates.

\paragraph{Intermediate periods require an additional input.}
If the spatial moment representations introduced below also hold,
then for even \(n_j\le L_x\le2n_j\) and \(m_j\le L_y\le2m_j\),
\[
	Z(L_x,L_y,\beta_j)
	\le Z(n_j,m_j,\beta_j)^{L_xL_y/(n_jm_j)}
	\le(1+r_j)^4.
\]
Since \((1+r)^4-1<5r\) for \(0<r\le0.01\),
\[
	e^{-\beta_j\Delta_{L_x,L_y}}
	\le R(L_x,L_y,\beta_j)<5r_j=20^{-2^j}.
\]
The same gap bound follows on these windows. In particular, equal seed
periods cover all sufficiently large even square tori. Scalar doubling
inequalities alone establish the dyadic result, not this interpolation.

\paragraph{Errors in spatial gluing.}
The exact inequalities can be relaxed. Suppose at \(b=2\beta_j\),
\begin{align*}
	Z(2n_j,m_j,b)  & \le e^{\eta_{x,j}}Z(n_j,m_j,b)^2,  \\
	Z(2n_j,2m_j,b) & \le e^{\eta_{y,j}}Z(2n_j,m_j,b)^2,
\end{align*}
with nonnegative errors. Put \(\eta_j=2\eta_{x,j}+\eta_{y,j}\).
The same proof gives the rigorous perturbed update
\begin{equation}
	\boxed{R_{j+1}\le e^{\eta_j}(1+R_j^2)^4-1.}
	\label{eq:ff-perturbed}
\end{equation}
For \(R_j,\eta_j\le0.01\), this is at most
\(5R_j^2+2\eta_j\). A constant nonzero gluing error is therefore
insufficient for this iteration to prove a gap; it creates an error
floor. An estimate \(\eta_j\le cR_j^2\) instead retains a quadratic
recursion with coefficient \(5+2c\). PEPS boundary estimates would
need an explicit bridge to these thermal partition-function errors
before they could be used as this input.

\section{Directional Purities and Spatial Transfer Operators}

Let \(H_{n,m}\) be a finite-dimensional Hamiltonian on a periodic
two-dimensional lattice, and put
\[
	Z(n,m,\beta)=\Tr e^{-\beta H_{n,m}},\qquad \beta>0.
\]
The integers \(n,m\) count spatial slices or fixed unit cells, not
necessarily individual spins. We use even periods throughout.
Write \(v=(n,m,\beta)\), with commuting doubling operations
\[
	D_xv=(2n,m,\beta),\quad D_yv=(n,2m,\beta),\quad
	D_\tau v=(n,m,2\beta).
\]
Define
\begin{align}
	P_x(v) & =\frac{Z(2n,m,\beta)}{Z(n,m,\beta)^2}, &
	P_y(v) & =\frac{Z(n,2m,\beta)}{Z(n,m,\beta)^2}, &
	P_\tau(v) & =\frac{Z(n,m,2\beta)}{Z(n,m,\beta)^2}.
	\label{eq:purities}
\end{align}
Here \(\tau\) denotes imaginary time. The thermal purity \(P_\tau(v)\) is always
\(\Tr\rho_v^2\) for \(\rho_v=e^{-\beta H_{n,m}}/Z(v)\).

\paragraph{Spatial-transfer hypothesis.}
For every transverse period and inverse temperature used below, assume
there exist self-adjoint operators \(T_x(m,\beta)\) and
\(T_y(n,\beta)\), independent of their respective longitudinal periods,
such that
\begin{equation}
	Z(n,m,\beta)=\tr T_x(m,\beta)^n
	=\tr T_y(n,\beta)^m.
	\label{eq:transfer}
\end{equation}
Assume the even powers in these identities are positive trace-class
operators. No uniform bound on the transfer-space dimension is needed.
The operators for different directions need not act on the same space
or commute. The identities must be exact; an unbounded truncation error
cannot be inserted into the argument.

For even \(n,m\), the trace-one operators
\[
	\sigma_x(v)=\frac{T_x(m,\beta)^n}{Z(v)},\qquad
	\sigma_y(v)=\frac{T_y(n,\beta)^m}{Z(v)},\qquad
	\sigma_\tau(v)=\rho_v
\]
are positive, and \(P_a(v)=\tr\sigma_a(v)^2\). Consequently,
\begin{equation}
	0<P_a(v)\le1,\qquad
	\sigma_a(D_av)=\frac{\sigma_a(v)^2}{P_a(v)},
	\qquad a\in\{x,y,\tau\}.
	\label{eq:square-states}
\end{equation}
An exact tensor-network representation alone does not imply
\eqref{eq:transfer}: spatial reflection, positivity, and the identification
of the same transfer operator at different lengths must be checked.

\section{Squaring Lemma and Spacetime Cube Coupling}

For any positive trace-one operator with eigenvalues \(p_i\), write
\(P=\sum_i p_i^2\). Its normalized square has purity \(P'\), with
\[
	1-P'=\frac{2\sum_{i<j}p_i^2p_j^2}{P^2}
	\le\frac{(1-P)^2}{2P^2}.
\]
Thus, with
\begin{equation}
	f(z)=\frac{z^2}{2(1-z)^2},\qquad 0\le z<1,
	\label{eq:f}
\end{equation}
we have \(1-P_a(D_av)\le f(1-P_a(v))\).

Every pair of directions also satisfies the exact identity
\begin{equation}
	P_a(v)^2P_b(D_av)=P_b(v)^2P_a(D_bv).
	\label{eq:face}
\end{equation}
This follows by cancellation of partition functions and requires no
transfer assumption. The same cancellation on the three-dimensional
cube gives, for any permutation \((a,b,c)\) of \((x,y,\tau)\),
\begin{equation}
	Q(v):=\frac{Z(D_xD_yD_\tau v)}{Z(v)^8}
	=P_a(v)^4P_b(D_av)^2P_c(D_bD_av).
	\label{eq:cube}
\end{equation}
The exponents \(4,2,1\) count how often the earlier ratios enter the
successive squarings. In particular, since every factor is at most one,
\begin{equation}
	P_c(D_aD_bv)\ge Q(v).
	\label{eq:last-edge}
\end{equation}
The cube identity supplies the coupling; positivity supplies this
inequality.

\section{Staggered Three-Component Scale-Doubling Iteration}

At the base point \(v\), assume
\begin{equation}
	P_x(v)\ge1-\ee_x,\qquad
	P_y(D_xv)\ge1-\ee_y,\qquad
	P_\tau(D_yD_xv)\ge1-\ee_\tau.
	\label{eq:seed}
\end{equation}
These are the three edges of the path \(x,y,\tau\) through one cube.
Put \(v'=D_xD_yD_\tau v\). The next-scale quantities occupy the same
staggered positions:
\begin{equation}
	P_x(v'),\qquad P_y(D_xv'),\qquad P_\tau(D_yD_xv').
	\label{eq:target}
\end{equation}
In coordinates, the input and output are
\[
	\begin{array}{c|c|c}
		     & \text{input}        & \text{output}         \\ \hline
		x    & P_x(n,m,\beta)      & P_x(2n,2m,2\beta)     \\
		y    & P_y(2n,m,\beta)     & P_y(4n,2m,2\beta)     \\
		\tau & P_\tau(2n,2m,\beta) & P_\tau(4n,4m,2\beta).
	\end{array}
\]

\begin{theorem}[Three-component update]
	Define successively
	\begin{align}
		u_x         & =(1-\ee_x)^4(1-\ee_y)^2(1-\ee_\tau),         &
		\delta_x    & =f(1-u_x),\label{eq:update-x}                  \\
		u_y         & =(1-\ee_y)^4(1-\ee_\tau)^2(1-\delta_x),      &
		\delta_y    & =f(1-u_y),\label{eq:update-y}                  \\
		u_\tau      & =(1-\ee_\tau)^4(1-\delta_x)^2(1-\delta_y),   &
		\delta_\tau & =f(1-u_\tau).\label{eq:update-t}
	\end{align}
	Provided the input defects lie in \([0,1)\) and
	\(\delta_x,\delta_y<1\), the output defects in
	\eqref{eq:target} are bounded by \(\delta_x,\delta_y,\delta_\tau\),
	respectively.
\end{theorem}

\begin{proof}
	First use the cube at \(v\). Its path \(x,y,\tau\) gives
	\(Q(v)\ge u_x\). Reorder the path to \(y,\tau,x\); by
	\eqref{eq:last-edge},
	\(P_x(D_yD_\tau v)\ge u_x\).
	Squaring in \(x\) yields \(P_x(v')\ge1-\delta_x\).

	Next use the cube based at \(D_xv\). Its path \(y,\tau,x\) uses two
	old bounds and the just-proved bound:
	\[
		Q(D_xv)=P_y(D_xv)^4P_\tau(D_yD_xv)^2P_x(v')\ge u_y.
	\]
	Reordering to \(x,\tau,y\) gives
	\(P_y(D_x^2D_\tau v)\ge u_y\).
	Squaring in \(y\) yields \(P_y(D_xv')\ge1-\delta_y\).

	Finally use the cube based at \(D_xD_yv\). Its path \(\tau,x,y\)
	gives
	\[
		Q(D_xD_yv)=P_\tau(D_xD_yv)^4P_x(v')^2P_y(D_xv')\ge u_\tau.
	\]
	Reordering to \(x,y,\tau\) gives
	\(P_\tau(D_x^2D_y^2v)\ge u_\tau\).
	Squaring in \(\tau\) yields \(P_\tau(D_xD_yv')\ge1-\delta_\tau\).
\end{proof}

\begin{theorem}[Closed scalar recursion]
	If the three input defects are at most \(\ee\le10^{-3}\), then all
	three output defects are at most
	\begin{equation}
		g(\ee)=f\bigl(1-(1-\ee)^7\bigr)
		=\frac12\bigl[(1-\ee)^{-7}-1\bigr]^2
		\le25\ee^2.
		\label{eq:scalar}
	\end{equation}
\end{theorem}

\begin{proof}
	Bernoulli's inequality and monotonicity of \(f\) give
	\[
		g(\ee)\le f(7\ee)
		=\frac{49\ee^2}{2(1-7\ee)^2}
		\le25\ee^2\le\ee.
	\]
	The numerical inequality is exact: at \(\ee=1/1000\),
	\(49{,}000{,}000<50\cdot993^2=49{,}302{,}450\).
	The first update has \(u_x\ge(1-\ee)^7\), so
	\(\delta_x\le g(\ee)\le\ee\).
	This implies \(u_y\ge(1-\ee)^7\), hence the same bound on
	\(\delta_y\); it then implies the corresponding bound on
	\(\delta_\tau\). Thus the recursion closes.
\end{proof}

\section{Uniform Spectral Gap from Certified Seeds}

Suppose three finite-scale certificates establish
\begin{equation}
	P_x(n_0,m_0,\beta_0),\quad
	P_y(2n_0,m_0,\beta_0),\quad
	P_\tau(2n_0,2m_0,\beta_0)\ \ge0.999.
	\label{eq:certificates}
\end{equation}
Set
\begin{equation}
	n_j=2^jn_0,\quad m_j=2^jm_0,\quad\beta_j=2^j\beta_0,
	\qquad \ee_j=\frac1{25}40^{-2^j}.
	\label{eq:scales}
\end{equation}
Then \(\ee_0=10^{-3}\) and \(\ee_{j+1}=25\ee_j^2\), so the
same three staggered purities at scale \(j\) are at least
\(1-\ee_j\).

We can also cover intermediate even periods. Choose even \(L_x,L_y\)
such that, for the same \(j\),
\begin{equation}
	n_j\le L_x\le2n_j,\qquad m_j\le L_y\le2m_j.
	\label{eq:window}
\end{equation}
For fixed transverse period and temperature, \eqref{eq:transfer}
expresses \(Z\) as an even spectral moment. Monotonicity of the
\(\ell^p\) norms, applied successively in \(x\) and \(y\), gives
\begin{align*}
	Z(L_x,L_y,\beta_j)
	 & \le Z(n_j,m_j,\beta_j)^{L_xL_y/(n_jm_j)},     \\
	Z(L_x,L_y,2\beta_j)
	 & \ge Z(2n_j,2m_j,2\beta_j)^{L_xL_y/(4n_jm_j)}.
\end{align*}
With \(\alpha=L_xL_y/(4n_jm_j)\in[1/4,1]\), this implies
\begin{align}
	P_\tau(L_x,L_y,\beta_j)
	 & \ge Q(n_j,m_j,\beta_j)^\alpha\notag         \\
	 & \ge(1-\ee_j)^{7\alpha}\ge1-7\ee_j>\tfrac12.
	\label{eq:interpolation}
\end{align}
If \(w_0\ge w_1\ge\cdots\) are the physical Gibbs weights, purity
is at most \(w_0\). Thus the ground state is unique and
\[
	e^{-\beta_j\Delta_{L_x,L_y}}
	=\frac{w_1}{w_0}
	\le\frac{7\ee_j}{1-7\ee_j}
	<40^{-2^j}.
\]
Consequently,
\begin{equation}
	\boxed{\Delta_{L_x,L_y}>\frac{\log40}{\beta_0}.}
	\label{eq:gap}
\end{equation}
This holds on the union of windows \eqref{eq:window}, conditional on
\eqref{eq:transfer} and \eqref{eq:certificates}. In particular, taking
\(n_0=m_0\) covers every even square torus of side at least \(n_0\).
It does not assert a bound for all aspect ratios or for open boundaries.
For a common seed defect \(0<\ee_0\le10^{-3}\), the same argument
replaces \(40\) by \((25\ee_0)^{-1}\).

\section{Application and Obstructions in the Hexagonal AKLT Model}

For the spin-\(3/2\) hexagonal AKLT Hamiltonian in projector
normalization,
\[
	H_{n,m}^{\mathrm{AKLT}}=\sum_{\langle u,v\rangle}
	P_{uv}^{(S_{\mathrm{tot}}=3)},
\]
choose fixed unit cells and two compatible periodic slice directions.
The frustration-free refinement requires:
\begin{enumerate}
	\item Prove the two scalar gluing inequalities
	      \eqref{eq:ff-x}--\eqref{eq:ff-y} on the required scales, or
	      control their errors through \eqref{eq:ff-perturbed}.
	      Self-adjoint transfers are one sufficient construction.
	\item Bound one finite-volume partition-function excess, or use a
	      certified finite-volume gap and dimension as in
	      \eqref{eq:ff-finite-gap}. Spatial moment interpolation is a
	      separate requirement if all intermediate periods are desired.
\end{enumerate}
No AKLT scalar gluing proof is supplied here. The result is therefore
a proved conditional gap criterion, not an AKLT gap proof.

\paragraph{An exact obstruction to the standard reflection argument.}
This spatial input cannot simply be borrowed from the Heisenberg
interaction. Consider one spin-\(3/2\) AKLT bond
\(h=P^{(3)}\) and the quadrupole
\[
	Q=(S^z)^2-\frac54I=\operatorname{diag}(1,-1,-1,1).
\]
Let \(\theta(A)=U\overline A U^\dagger\), where
\(U=e^{-i\pi S^y}\), be the usual spin time-reversal action used with
a spatial reflection. Then \(\theta(Q)=Q\). Reflection positivity
for this bond would require
\(\Tr[e^{-t h}(A\otimes\theta(A))]\ge0\) for every \(A\).
But \(\Tr Q=0\), \(e^{-tP^{(3)}}=I+(e^{-t}-1)P^{(3)}\), and
\begin{equation}
	\Tr[P^{(3)}(Q\otimes Q)]=\frac75.
	\label{eq:quadrupole-trace}
\end{equation}
Hence for every \(t>0\),
\begin{equation}
	\boxed{\Tr[e^{-tP^{(3)}}(Q\otimes\theta(Q))]
		=\frac75(e^{-t}-1)<0.}
	\label{eq:rp-obstruction}
\end{equation}
For a direct check of \eqref{eq:quadrupole-trace}, the diagonal
expectations of \(Q\otimes Q\) in \(\lvert3,M\rangle\),
\(M=-3,-2,\ldots,3\), are
\(1,-1,1/5,1,1/5,-1,1\); their sum is \(7/5\).
Equivalently, with \(X=\mathbf S_1\cdot\mathbf S_2\), the projector is
\(P^{(3)}=(X+15I/4)(X+11I/4)(X+3I/4)/90\).
The characterization of thermal reflection positivity by cross-plane
couplings is discussed in Ref.~\cite{jaffe}.

This calculation rules out this particular local reflection-positive
construction. It does not disprove the weaker scalar inequalities
\eqref{eq:ff-x}--\eqref{eq:ff-y}, or every possible blocked transfer
construction. Those statements must be tested independently.

There is also a direct spectral diagnostic for scalar gluing. On two
tori with unique zero-energy ground states, write
\(Z(n,m,\beta)=1+g_{n,m}e^{-\beta\Delta_{n,m}}+\cdots\).
If \(\Delta_{2n,m}<\Delta_{n,m}\), then
\(Z(2n,m,\beta)>Z(n,m,\beta)^2\) for all sufficiently large
\(\beta\). Thus a verified decrease of the finite-volume gap under
doubling would obstruct the exact gluing inequality at those sizes.

The PEPS boundary method offers a complementary route: relative
factorization of overlapping ground-state boundary operators implies
control of regional ground-space projectors and hence a bulk
gap~\cite{klp}. The recent AKLT preprint~\cite{kastoryano} develops
boundary-response estimates using paths and loops. Such estimates may
inform the model-specific analysis above, but they are not automatically
thermal-purity certificates. Here the objects are the full physical
partition functions, rather than the PEPS boundary Gram operators.

\section{Empirical Doubling Coefficients and Diagnostics}

Exact spatial gluing is sufficient but not necessary for the recursion.
For the dyadic sequence define the actual coefficient
\begin{equation}
	C_j=\frac{R(n_{j+1},m_{j+1},\beta_{j+1})}
	{R(n_j,m_j,\beta_j)^2}.
	\label{eq:actual-C}
\end{equation}
If \(C_j\le C\) uniformly, with \(C\ge1\) and \(CR_0<1\), then
induction gives
\[
	R_j\le C^{-1}(CR_0)^{2^j},\qquad
	\Delta_{n_j,m_j}\ge\frac{\log(1/(CR_0))}{\beta_0}.
\]
This only needs control along the chosen sequence
\(\beta_j/n_j=\beta_0/n_0\), rather than at all temperatures for
each fixed torus. Finitely many measured coefficients do not prove the
uniform bound.

A useful diagnostic of the sufficient gluing construction is
\[
	G_j=\log Z(2n_j,2m_j,2\beta_j)
	-4\log Z(n_j,m_j,2\beta_j).
\]
Taking \(\eta_j=\max\{0,G_j\}\) in
\eqref{eq:ff-perturbed} is valid directly, and can be sharper than
bounding the two directional errors separately.

\subsection{First computational pilot}

The companion code in \texttt{models/aklt/} uses primitive cells
\(A(x,y),B(x,y)\), with edges from \(A(x,y)\) to
\(B(x,y),B(x-1,y),B(x,y-1)\) and both coordinates periodic.
The initial geometries are \(2\times2\), \(4\times2\), and
\(4\times4\) cells, containing 8, 16, and 32 spins. These small
quotients have finite-size winding effects; in particular, the
eight-spin quotient has short winding cycles. They are not open
hexagonal patches.

Independent full-spectrum diagonalization of the eight-spin quotient
includes all \(4^8=65536\) states through total-\(S^z\) sectors.
It gives one eigenvalue within \(10^{-9}\) of zero and
\begin{align*}
	E_0       & \simeq1.35\times10^{-14}, &
	\Delta_8  & \simeq0.0920298633241,      \\
	R(2,2,64) & \simeq0.00831113960.      &
\end{align*}
The numerical threshold for \(R\le0.01\) is
\(\beta\simeq61.9932\). This provides a candidate finite seed at
\(\beta_0=64\), not an interval certificate. If the seed and uniform
spatial gluing were certified, the conservative recursion above would
give \((\log20)/64\simeq0.0468083\). Neither the full gap claim nor
uniform gluing follows from this finite calculation.

The thermal implementation evolves the MPO amplitude
\(e^{-\beta H/2}\), using
\(e^{-tP^{(3)}}=I+(e^{-t}-1)P^{(3)}\) for each local gate.
It retains the full thermal trace and reconstructs \(\log Z\) from
normalization factors. A forward/reverse gate sweep has second-order
Trotter error; compression and the swaps needed for nonlocal gates
introduce additional errors.

At \(\beta=2\), full diagonalization gives
\(\log Z=5.31205522615\). The following benchmark uses time step
\(\delta\beta=1/8\), cutoff \(10^{-12}\), and the default policy
of returning moved sites after each gate:
\begin{center}
	\begin{tabular}{rrr}
		\toprule
		Maximum MPO bond dimension & Computed \(\log Z\) & Error in \(\log Z\) \\
		\midrule
		16                         & 3.53951535          & -1.77253988         \\
		32                         & 3.88853480          & -1.42352043         \\
		64                         & 4.72481476          & -0.58724047         \\
		128                        & 4.97859393          & -0.33346129         \\
		\bottomrule
	\end{tabular}
\end{center}
Even this moderate-temperature benchmark is not converged. For example,
at bond dimension 64, halving the time step gives
\(\log Z=4.58083645\), increasing the discrepancy. The smaller time
step alone does not remove compression error. Negative gluing defects
or small doubling coefficients from these compressed calculations must
therefore not be treated as evidence for the desired inequalities.
A separate test keeps moved sites between gates and restores their
order only after each complete step. At bond dimension 128 and
\(\delta\beta=1/8\), this gives \(\log Z(2)=4.91941052\);
reducing repeated swaps does not resolve the benchmark discrepancy.

Machine-readable spectra, convergence grids, and diagnostics are saved
under \path{models/aklt/results/}, with source hashes, geometry, parameters,
completion flags, and \texttt{certified=false}. The smallest-system
checks compare the two independent projector constructions, the exact
single-bond partition function, and full-spectrum four-site evolution,
including its nonlocal closing bond. The immediate computational task
is to resolve partition-function differences accurately before using
larger thermal calculations to assess the recursion.

\section{Quadratic concentration without a self-adjoint spatial transfer}
\label{sec:nonreal}

The coupled squaring mechanism of the OpenAI proof~\cite{oai,main} uses
scalar properties of the spatial moments. Self-adjointness is one
sufficient way to obtain those properties. Nonreal transfer eigenvalues
rule out a positive metric making that transfer self-adjoint, but do not
by themselves rule out the scalar inequalities needed for the recursion.
We establish this distinction for one exact AKLT projection-filter slice.

\subsection{An exact filter and its spatial transfer}

Let \(E_1,E_2,E_3\) be the three matchings used in the pilot, with
\(E_1\) intracell, \(E_2\) joining \(A(x,y)\) to \(B(x-1,y)\),
and \(E_3\) joining \(A(x,y)\) to \(B(x,y-1)\). Set
\[
	q_e=I-P_e^{(3)},\qquad Q_c=\prod_{e\in E_c}q_e,\qquad
	D=Q_3Q_2Q_1,\qquad F=D^\dagger D=Q_1Q_2Q_3Q_2Q_1.
\]
Then \(0\le F\le I\) and \(\ker(I-F)=\ker H\): equality of
norms through a product of orthogonal projections requires every
projection to fix the vector. Define the discrete-depth moments
\[
	\mathcal Z(n,m,t)=\Tr F_{n,m}^{\,t},\qquad t\in\mathbb N.
\]
Temporal squaring is exact for these moments. The converse detectability
inequality gives \(I-F\le4H\)~\cite{aav}. Thus a uniform estimate
\(F\le G+q(I-G)\), where \(G\) projects onto \(\ker H\) and
\(q<1\), would give \(\Delta(H)\ge(1-q)/4\).

Fix transverse width \(n=2\) and depth \(t=1\), and let the
longitudinal period be \(L\). A row contains the four spins
\(A_0,A_1,B_0,B_1\). Write
\[
	A=Q_{1,\mathrm{row}},\qquad
	B=Q_{2,\mathrm{row}}A Q_{2,\mathrm{row}},\qquad
	\mathcal R(X)_{(a,a'),(b,b')}=X_{(a,b),(a',b')}.
\]
Here each of \(a,b\) indexes two spins, hence has dimension 16.
Cyclicity, \(Q_1^2=Q_1\), and direct contraction of the interrow
\(Q_3\) bonds give the \(256\)-dimensional transfer
\begin{equation}
	T=\mathcal R(A)^{\mathsf T}\mathcal R(B^{\mathsf T}),\qquad
	z_L:=\mathcal Z(2,L,1)=\tr T^L.
	\label{eq:filter-row}
\end{equation}
This transfer has a conjugate pair
\(0.4111133913\pm0.3495535140i\), each with multiplicity five.
Exact characteristic-polynomial and root-count checks establish their
nonreality. Therefore neither \(T\) nor \(T^2\) can be made
self-adjoint by a positive metric. The scalar moments nevertheless
satisfy a quadratic concentration estimate.

\subsection{A scalar replacement for spectral positivity}

\begin{lemma}[Concentration from a controlled trace remainder]
	Suppose \(z_L=\lambda_0^L(1+r_L)\), where \(\lambda_0>0\), and
	for every even \(L\ge L_0\),
	\[
		0<r_L\le\tfrac1{10},\qquad
		0<r_{2L}\le\kappa r_L^2,\qquad 0<\kappa\le\tfrac{34}{100}.
	\]
	Then the scalar defect \(d_L=1-z_{2L}/z_L^2\) obeys
	\begin{equation}
		0<d_L<1,\qquad
		d_{2L}\le\frac{\kappa}{2}\left(\frac{11}{10}\right)^4 d_L^2
		<\frac14d_L^2.
		\label{eq:nonreal-square}
	\end{equation}
\end{lemma}

\begin{proof}
	The identity
	\[
		d_L=\frac{2r_L+r_L^2-r_{2L}}{(1+r_L)^2}
	\]
	gives \(d_L\ge2r_L/(1+r_L)^2>0\), since \(\kappa<1\),
	and \(d_L<1\) since \(z_{2L}>0\). Also
	\(d_{2L}\le2r_{2L}\le2\kappa r_L^2\), using
	\(r_{4L}\ge0\). Substitute
	\(r_L\le(1+r_L)^2d_L/2\) and \(r_L\le1/10\).
\end{proof}

\paragraph{Exact spectral bounds for the fixed slice.}
The verifier finds one simple leading eigenvalue \(\lambda_0\),
one real eigenvalue \(\alpha\) of multiplicity three, 242 remaining
real eigenvalues, and ten nonreal eigenvalues. Exact rational root
counts certify
\begin{align*}
	9.57      & <\lambda_0<9.59,                               & 3.36 & <\alpha<3.37, \\
	|\lambda| & <0.84\quad\text{for the remaining real modes}, &      &               \\
	|\lambda| & <0.55\quad\text{for every nonreal mode}.       &      &
\end{align*}
All displayed decimal endpoints are exact rational bounds. No
diagonalizability assumption is needed: traces of powers count
eigenvalues with algebraic multiplicity.

For even \(L\ge4\), let
\(r_L=\sum_{i>0}(\lambda_i/\lambda_0)^L\) and
\(u=55/336\). Real modes contribute nonnegatively, so
\[
	r_L\ge3(\alpha/\lambda_0)^L
	\left[1-\frac{10}{3}u^L\right]>0.
\]
The absolute remainder has the uniform bound
\[
	r_L\le3\left(\frac{337}{957}\right)^4
	+242\left(\frac{84}{957}\right)^4
	+10\left(\frac{55}{957}\right)^4
	<0.060605<\frac1{10}.
\]
Bounding the numerator at \(2L\) and the denominator at \(L\) gives
\[
	\frac{r_{2L}}{r_L^2}
	\le
	\frac{3+242(1/4)^{2L}+10u^{2L}}
	{9[1-(10/3)u^L]^2}
	\le\kappa_*:=
	\frac{3+242(1/4)^8+10u^8}
	{9[1-(10/3)u^4]^2}
	<0.335348<0.34.
\]
The lemma therefore proves
\begin{equation}
	\boxed{d_{2L}\le\tfrac14d_L^2
		\quad\text{for every even }L\ge4,
		\quad z_L=\mathcal Z(2,L,1).}
	\label{eq:fixed-slice-result}
\end{equation}
This is stronger than the scalar squaring inequality
\(d_{2L}\le f(d_L)\) used earlier, although the spatial transfer
has nonreal eigenvalues.

\paragraph{Verification and remaining scope.}
In the physical basis change \(S=\operatorname{diag}\sqrt{\binom3a}\),
\(a=0,1,2,3\), each bond gate becomes rational with denominator 60, so
\(60^8T\) is integral in the corresponding transfer basis; exact
characteristic polynomials and root counts establish the spectral bounds above
and match an independent eight-spin tensor contraction. The Julia pilot in
\path{models/aklt/} (\path{scripts/run.jl}, \path{test/runtests.jl}) checks the
bond projector polynomial, the quadrupole reflection obstruction
\eqref{eq:quadrupole-trace}, and the exact and thermal-purification benchmarks.

Equation~\eqref{eq:fixed-slice-result} holds at one fixed transverse
width and one fixed depth. It does not control the mixed-width,
mixed-depth tuples needed by the cube recursion. The next mathematical
target is to propagate such domination estimates in both transverse
size and depth, so that both spatial ratios remain in \((0,1]\) and
obey their squaring bounds throughout the staggered recursion.
The cube proof then applies to the scalar ratios directly, without
requiring spatial density operators. No uniform AKLT gap follows from
the fixed-slice certificate alone.

\begin{thebibliography}{9}
	\setlength{\itemsep}{2pt plus 1pt}
	\setlength{\parskip}{0pt}
	\bibitem{klp}
	M. J. Kastoryano, A. Lucia, and D. P\'erez-Garc\'ia,
	\emph{Locality at the boundary implies gap in the bulk for 2D PEPS},
	Commun. Math. Phys. \textbf{366}, 895--926 (2019),
	\href{https://arxiv.org/abs/1709.07691}{arXiv:1709.07691}.
	\bibitem{kastoryano}
	M. J. Kastoryano,
	\emph{Spectral Gap of the Hexagonal AKLT Model via Boundary-State
		Factorization}, preprint (September 23, 2026),
	\href{https://arxiv.org/abs/2609.27993}{arXiv:2609.27993}.
	\bibitem{jaffe}
	A. Jaffe and B. Janssens,
	\emph{Characterization of Reflection Positivity: Majoranas and Spins},
	\href{https://arxiv.org/abs/1506.04197}{arXiv:1506.04197}.
	\bibitem{oai}
	OpenAI, \emph{The periodic spin-one Haldane gap},
	manuscript dated September 24, 2026,
	\href{https://github.com/openai/math/blob/main/preprints/The-periodic-spin-one-Haldane-gap-September-24-2026/paper.pdf}{OpenAI Math Release}.
	\bibitem{main}
	Companion note, \emph{A diagrammatic proof of the Haldane gap via spacetime
		tensor networks}, \texttt{notes/thermal-mpo-gap-note.tex} (2026).
	\bibitem{aav}
	A. Anshu, I. Arad, and T. Vidick,
	\emph{A simple proof of the detectability lemma and spectral gap amplification},
	\href{https://arxiv.org/abs/1602.01210}{arXiv:1602.01210}.
\end{thebibliography}
\end{document}
